\chapter{Dane}

% opis kształtu i znaczenia danych
% 1.dane (co mierzymy, jak to wygląda, co uzyskujemy)
\section{Obiekt pomiaru}
% mierzymy wartości z tensometru przetwornikami ADC kamizelką
% wygląd danych - naciąg każdego z pasów
% uzyskujemy - linie jsonl z wartościami kolejnych próbek

% 2.transformacje danych (przejscia pomiedzy etapami) - Flow
\section{Etapy przetwarzania danych}
% etapy:
% I.handle input data, MJ
% II.breath separation, JD
\subsection{Oddzielanie oddechów}
Celem tego etapu jest wstępne podzielenie sygnału na mniejsze części reprezentujące
kolejne oddechy w analizowaynm pomiarze. Realizowane jest to poprzez wyznaczenie 
wszystkich maksymów i minimów lokalnych oraz zapisanie ich w programie tak, by
przygotować je do dalszej obróbki.

\begin{itemize}
    \item \textbf{Dane wejściowe:} \\
    Zbiór wartości odstawania od średniej wartości sygnału na danym pasie, dla każdego
    pasa w kamizelce pomiarowej. Każda z tych wartości przypisany ma odpowiedni 
    znacznik czasu od rozpoczęcia pomiaru, w którym pobrano próbkę.
    \item \textbf{Dane wyjściowe:} \\
    Zbiór punktów reprezentujących wszystkie maksyma i minima lokalne sygnału 
    wraz z przypisanymi im wartościami znaczników czasu.
\end{itemize}

% III.outlier detection, MJ
% IV.split data into segments, MJ
% V.feature extraction, JD
\subsection{Ekstrakcja cech}
Po podzieleniu danych na segmenty, a więc ostatecznym przygotowaniu ich 
do analizy, należy rozpocząć proces pozyskiwania informacji potrzebnych do 
prawidłowego działania klasyfikatora. Sama sieć neuronowa nie potrzebuje 
zdefiniowanych cech sygnału, by funkcjonować, jednak w celach badawczych, 
zdecydowano się na ich ekstrakcję ze względu na możliwość sprawdzenia różnic
we wzorcach oddechowych objętych analizą osób.

\begin{itemize}
    \item \textbf{Dane wejściowe:} \\
    Zbiór plików zawierających podzielone na segmenty dane z pomiarów po 
    wszystkich przekształceniach, w których nazwa wyznacza etykietę, ponieważ 
    zawiera typ aktywności oraz identyfikator osoby
    \item \textbf{Dane wyjściowe:} \\
    Plik wyjściowy
    \texttt{extracted\_features.jsonl}
    w formacie
    \texttt{.jsonl}, zawierający wektory 
    cech wyznaczone dla kolejnych segmentów danych pomiarowych. 
\end{itemize}

\noindent\textbf{Cechy składające się na wektor dla segmentu:}

\begin{itemize}
    \item \texttt{bpm} - liczba oddechów na minutę
    \item \texttt{breath\_depth} - wyznaczone poprzez uśrednienie wartości szczytów 
    oddechów w całym segmencie dla środkowego pasa
    \item \texttt{breath\_depth\_std} - odchylenie standardowe sygnału, jako głębokość 
    oddechu dla środkowego pasa
    \item \texttt{belt\_share} - udział poszczególnych pasów w oddychaniu, reprezentowany 
    przez 5 wartości, po jednej dla każdego pasa. Jest on wyznaczany poprzez 
    zsumowanie wartości głębokości oddechu dla każdego pasa, tworzy to mianownik 
    ułamka, który liczony jest osobno dla każdego tensometru, gdzie licznikiem jest 
    głębokość oddechu dla tego konkretnego sensora
    \item \texttt{belt\_share\_std} - wyznaczany jest on analogicznie do powyższego, korzysta 
    jednak z wyżej wspomnianych odchyleń standardowych sygnału, jako głębokości
    \item \texttt{breathing\_phase\_lengths} - jest to wartość średnia czasu trwania kolejnych 
    faz oddychania, a mianowicie: faza wdechu, pauza po wdechu, faza wydechu, 
    pauza po wydechu
\end{itemize}

% VI. Visualization (separate) JD
\section{Wizualizacja danych}